\documentclass[11pt,a4paper]{article}

% --- Pacchetti ---
\usepackage[utf8]{inputenc}
\usepackage[T1]{fontenc}
\usepackage{lmodern}
\usepackage[italian]{babel}
\usepackage[margin=1in]{geometry}
\usepackage{xcolor}
\usepackage{listings}
\usepackage{hyperref}
\usepackage{enumitem}

% --- Colori ---
\definecolor{primary}{RGB}{38, 66, 139}
\definecolor{codebg}{RGB}{245, 245, 245}

\hypersetup{colorlinks=true, linkcolor=primary, urlcolor=primary}

% --- Configurazione Codice ---
\lstdefinelanguage{javascript}{
    keywords={async, await, break, case, catch, class, const, continue, debugger, default, delete, do, else, export, extends, false, finally, for, function, if, import, in, instanceof, new, null, return, super, switch, this, throw, true, try, typeof, var, void, while, with, yield, let},
    morekeywords={console, log, Error, fetch, JSON, parse, stringify, useState, useEffect},
    sensitive=true,
    morecomment=[l]{//},
    morecomment=[s]{/*}{*/},
    morestring=[b]',
    morestring=[b]",
    morestring=[b]`,
    keywordstyle=\color{primary}\bfseries,
    commentstyle=\color{green!50!black}\itshape,
    stringstyle=\color{red!70!black},
    basicstyle=\ttfamily\small
}

\lstset{
    backgroundcolor=\color{codebg},
    basicstyle=\ttfamily\small,
    breaklines=true,
    frame=single,
    language=javascript
}

\title{\textbf{\color{primary} Esercitazione: Todo App con Next.js e React}}
\author{Programmazione e Tecnologie Web - Lezione 5}
\date{A.A. 2025/2026}

\begin{document}

\maketitle

\section{Obiettivo}
L'obiettivo di questa esercitazione è riscrivere il frontend della Todo App utilizzando **Next.js**. Passeremo dalla manipolazione diretta del DOM ad un approccio dichiarativo basato su componenti e stato reattivo.

\section{Esercizio 1: Setup del Progetto}
Utilizzeremo la CLI di Next.js per inizializzare l'ambiente.

\textbf{Task Live}:
\begin{enumerate}
    \item Aprire il terminale e digitare:
    \begin{lstlisting}[language=bash]
npx create-next-app@latest my-todo-react
    \end{lstlisting}
    \item Selezionare le opzioni: App Router (Yes), Tailwind CSS (Yes), TypeScript (No).
    \item Avviare il server di sviluppo con \texttt{npm run dev} e verificare che la pagina "Hello World" sia visibile su \texttt{http://localhost:3000}.
\end{enumerate}

\section{Esercizio 2: Componente Reattivo e Stato Locale}
Nel file \texttt{app/page.js}, implementate una lista di task utilizzando dati iniziali hardcoded.

\begin{lstlisting}[language=javascript]
"use client";
import { useState } from "react";

export default function Home() {
  const [tasks, setTasks] = useState([
    { id: 1, text: "Imparare React", done: false },
    { id: 2, text: "Configurare Next.js", done: true }
  ]);

  return (
    <main className="p-10">
      <h1 className="text-2xl font-bold">Todo List</h1>
      {/* Mappare i task qui */}
    </main>
  );
}
\end{lstlisting}

\section{Esercizio 3: Ricerca e Filtraggio (Stato Derivato)}
Aggiungete una barra di ricerca. La lista dei task visualizzati deve aggiornarsi in tempo reale mentre l'utente scrive.

\textbf{Task}:
\begin{enumerate}
    \item Creare una variabile di stato \texttt{search} per gestire l'input.
    \item Creare una variabile \texttt{filteredTasks} calcolata durante il render (non uno state!).
    \item Utilizzare \texttt{.filter()} e \texttt{.map()} per visualizzare i risultati.
\end{enumerate}

\section{Esercizio 4: Connessione al Backend Reale}
Utilizzate l'hook \texttt{useEffect} per scaricare i task dal server backend della Lezione 4.

\textbf{Nota}: Ricordate di abilitare il CORS nel server Flask o Express, oppure utilizzate le API Routes di Next.js per fare da proxy.

\section{Task Bonus (Avanzato)}
Implementate una **API Route** in \texttt{app/api/stats/route.js} che restituisca un JSON con il numero di task completati e quelli ancora da fare. Consumate questa API nel frontend.

\end{document}
